% Faithful formal-publication transcription of Section4.1, Lemma8 and entire proof. Published graphical arrays are preserved as vector crops.
\section*{4 Search-to-decision reduction by restricting to monomial groups}
\subsection*{4.1 The gadget restricting to the monomial group}
In this section, we present the gadget that restricts to the monomial group in the setting of Alternating Matrix Space Isometry. To show this, we will need the concept of monomial isometry; see Some Groups above. Recall that a matrix is monomial if, equivalently, it can be written as $DP$ where $D$ is a nonsingular diagonal matrix and $P$ is a permutation matrix. We say two matrix spaces $\mathcal A,\mathcal B$ are monomially isometric if there is some $M\in\operatorname{Mon}(n,\mathbb F)$ such that $M^t\mathcal A M=\mathcal B$.
\begin{authorclaim}{Lemma8}
Alternating Matrix Space Monomial Isometry reduces to Alternating Matrix Space Isometry. More specifically, there is a $\operatorname{poly}(n,m)$-time algorithm $r$ taking alternating matrix tuples to alternating matrix tuples, such that for $\mathbf A,\mathbf B\in\Lambda(n,\mathbb F)^m$, the matrix spaces $\mathcal A=\langle\mathbf A\rangle$ and $\mathcal B=\langle\mathbf B\rangle$ are monomially isometric if and only if the matrix spaces $\langle r(\mathbf A)\rangle$ and $\langle r(\mathbf B)\rangle$ are isometric.
\end{authorclaim}
The gadget used in Lemma8 is essentially to apply the gadget in Proposition7 ``in two directions.'' Still, to prove the correctness requires some work.
\begin{proof}
For $\mathbf A=(A_1,\ldots,A_m)\in\Lambda(n,\mathbb F)^m$, define $r(\mathbf A)$ to be the alternating matrix tuple $\widetilde{\mathbf A}=(\widetilde A_1,\ldots,\widetilde A_{m+n^2})\in\Lambda(n+n^2,\mathbb F)^{m+n^2}$, where
\begin{enumerate}
\item For $k=1,\ldots,m$, $\widetilde A_k=\begin{bmatrix}A_k&0\\0&0\end{bmatrix}$.
\item For $k=m+(i-1)n+j$, $i\in[n],j\in[n]$, $\widetilde A_k$ is the elementary alternating matrix $E_{i,in+j}-E_{in+j,i}$.
\end{enumerate}
At this point, some readers may wish to look at the large matrix in Equation2 and/or at Figure2.
It is clear that $r$ can be computed in time $\widetilde O((m+n^2)(n^2+n))=\operatorname{poly}(n,m)$. Given alternating matrix tuples $\mathbf A,\mathbf B$, let $\mathcal A,\mathcal B$ be the corresponding matrix spaces they span, and let $\widetilde{\mathcal A}=\langle r(\mathbf A)\rangle$ and $\widetilde{\mathcal B}=\langle r(\mathbf B)\rangle$. We claim that $\mathcal A$ and $\mathcal B$ are monomially isometric if and only if $\widetilde{\mathcal A}$ and $\widetilde{\mathcal B}$ are isometric.
To prove this, it will help to think of our matrix tuples $\mathbf A,\widetilde{\mathbf A}$, etc. as (corresponding to) 3-way arrays, and to view these 3-way arrays from two different directions. Towards this end, write the 3-way array corresponding to $\mathbf A$ as
\[\mathtt A=\begin{bmatrix}0&a_{1,2}&a_{1,3}&\cdots&a_{1,n}\\-a_{1,2}&0&a_{2,3}&\cdots&a_{2,n}\\-a_{1,3}&-a_{2,3}&0&\cdots&a_{3,n}\\\vdots&\vdots&\vdots&\ddots&\vdots\\-a_{1,n}&-a_{2,n}&-a_{3,n}&\cdots&0\end{bmatrix},\]
where $a_{i,j}$ are vectors in $\mathbb F^m$ (``coming out of the page''), namely $a_{i,j}(k)=A_k(i,j)$. The frontal slices of this array are precisely the matrices $A_1,\ldots,A_m$.
The 3-way array corresponding to $\widetilde{\mathbf A}=r(\mathbf A)$ is then the $(n+1)n\times(n+1)n\times(m+n^2)$ array:
\begin{center}\includegraphics[width=.98\linewidth]{sources/equation2.pdf}\end{center}
where $\widetilde a_{i,j}=\begin{bmatrix}a_{i,j}\\0\end{bmatrix}\in\mathbb F^{m+n^2}$ (here think of the vector $a_{i,j}$ as a column vector, not coming out of the page; in the above array we then lay the column vector $\widetilde a_{i,j}$ ``on its side'' so that it is coming out of the page), and $e_{i,j}:=e_{m+(i-1)n+j}\in\mathbb F^{m+n^2}$, which we can equivalently write as $\begin{bmatrix}0^m\\e_i\otimes e_j\end{bmatrix}$, where we think of $e_i\otimes e_j$ here as a vector of length $n^2$. Note that all the the nonzero blocks besides upper-left ``$\mathtt A$'' block only have nonzero entries that are strictly behind the nonzero entries in the upper-left block.
\begin{center}\includegraphics[width=.5\linewidth]{sources/figure2.pdf}\end{center}
\noindent Figure2. Pictorial representation of the reduction for Lemma8.
The second viewpoint, which we will also use below, is to consider the lateral slices of $\widetilde{\mathtt A}$, or equivalently, to view $\widetilde{\mathtt A}$ from the side. When viewing $\widetilde{\mathtt A}$ from the side, we see the $(n+1)n\times(m+n^2)\times(n+1)n$ 3-way array:
\begin{center}\includegraphics[width=.98\linewidth]{sources/equation3.pdf}\end{center}
where every $\ell_{i,k}\in\mathbb F^{n^2+n}$ has only the first $n$ components being possibly non-zero, namely, $\ell_{i,k}(j)=A_k(i,j)$ for $i\in[n],j\in[n],k\in[m]$ and $\ell_{i,k}(j)=0$ for any $j>n$.
\paragraph{For the only if direction.}
Suppose there exist $P\in\operatorname{Mon}(n,\mathbb F)$ and $Q\in\mathrm{GL}(m,\mathbb F)$, such that $P^t\mathbf A P=\mathbf B^Q$. We can construct $\widetilde P\in\operatorname{Mon}(n+n^2,\mathbb F)$ and $\widetilde Q\in\mathrm{GL}(m+n^2,\mathbb F)$ such that $\widetilde P^t\widetilde{\mathbf A}\widetilde P=\widetilde{\mathbf B}^{\widetilde Q}$. In fact, we will show that we can take $\widetilde P=\begin{bmatrix}P&0\\0&P'\end{bmatrix}$ where $P'\in\operatorname{Mon}(n^2,\mathbb F)$, and $\widetilde Q=\begin{bmatrix}Q&0\\0&Q'\end{bmatrix}$ where $Q'\in\operatorname{Mon}(n^2,\mathbb F)$. It is not hard to see that this form already ensures that the first $m$ matrices in the vector $\widetilde P^t\widetilde{\mathbf A}\widetilde P$ and those of $\widetilde{\mathbf B}^{\widetilde Q}$ are the same, since when $\widetilde P,\widetilde Q$ are of this form, those first $m$ matrices are controlled entirely by the $P$ (resp., $Q$) in the upper-left block of $\widetilde P$ (resp., $\widetilde Q$).
The remaining question is then how to design appropriate $P'$ and $Q'$ to take care of the last $n^2$ matrices in these tuples. This actually boils down to applying the following simple identity, but ``in 3 dimensions:'' Let $P$ be the permutation matrix corresponding to $\sigma\in S_n$, so that $Pe_i=e_{\sigma(i)}$, and $e_i^tP=e_{\sigma^{-1}(i)}^t$. Let $D=\operatorname{diag}(\alpha_1,\ldots,\alpha_n)$ be a diagonal matrix. Then
\begin{equation}\tag{4}P^tDP=\operatorname{diag}(\alpha_{\sigma^{-1}(1)},\ldots,\alpha_{\sigma^{-1}(n)}).\end{equation}
To see how Equation4 helps in our setting, it is easier to focus attention on the lower right $n^2\times n^2$ sub-array of $\mathtt A^{\mathrm{lat}}$, which can be represented as a symbolic matrix
\[M=\begin{bmatrix}x_1I_n&0&\cdots&0\\0&x_2I_n&\cdots&0\\\vdots&\vdots&\ddots&\vdots\\0&0&\cdots&x_nI_n\end{bmatrix}.\]
Here we think of the $x_i$'s as independent variables, whose indices correspond to ``how far into the page'' they are. That is, $x_i$ corresponds to the vector $\vec e_i$ in $\mathtt A^{\mathrm{lat}}$, which is coming out of the page and has its only nonzero entry $i$ slices back from the page.
Then the action of $P$ permutes the $x_i$'s and multiplies them by some scalars, the action of $P'$ is on the left-hand side, and the action of $Q'$ is on the right-hand side. Let $\sigma$ be the permutation supporting $P$. Then $P$ sends $M$ to
\[M^P=\begin{bmatrix}\alpha_{\sigma(1)}x_{\sigma(1)}I_n&0&\cdots&0\\0&\alpha_{\sigma(2)}x_{\sigma(2)}I_n&\cdots&0\\\vdots&\vdots&\ddots&\vdots\\0&0&\cdots&\alpha_{\sigma(n)}x_{\sigma(n)}I_n\end{bmatrix}.\]
So setting $P'=\sigma\otimes I_n$, $Q'$ the monomial matrix supported by $\sigma\otimes I_n$ with scalars being $1/\alpha_i$'s, we have $P'^tM^PQ'=M$ by Equation4.
\paragraph{For the if direction.}
Suppose there exist $\widetilde P\in\mathrm{GL}(n+n^2,\mathbb F)$ and $\widetilde Q\in\mathrm{GL}(m+n^2,\mathbb F)$, such that $\widetilde P^t\widetilde{\mathbf A}\widetilde P=\widetilde{\mathbf B}^{\widetilde Q}$. The key feature of these gadgets now comes into play: consider the lateral slices of $\widetilde{\mathtt A}$, which are the frontal slices of $\mathtt A^{\mathrm{lat}}$ (which may be easier to visualize by looking at Equation3 and Figure2). The first $n$ lateral slices of $\widetilde{\mathtt A}$ and $\widetilde{\mathtt B}$ are of rank $\ge n$ and $<2n$, while the other lateral slices are of rank $<n$ (in fact, they are of rank 1; note that without loss of generality we may assume $n>1$, for the only $1\times1$ alternating matrix space is the zero space). Furthermore, left multiplying a lateral slice by $\widetilde P^t$ and right multiplying it by $\widetilde Q$ does not change its rank. However, the action of $\widetilde P$ here is by $\widetilde P^t\widetilde{\mathbf A}\widetilde P$, and while the $\widetilde P^t$ here corresponds to left multiplication on the lateral slices (=frontal slices of $\mathtt A^{\mathrm{lat}}$), the $\widetilde P$ on the right here corresponds to taking linear combinations of the lateral slices. In other words, just as $\mathtt A^{\mathrm{lat}}$ is the ``side view'' of $\widetilde{\mathtt A}$, $(\widetilde P^t\mathtt A^{\mathrm{lat}}\widetilde Q)^{\widetilde P}$ is the side view of $(\widetilde P^t\widetilde{\mathtt A}\widetilde P)^{\widetilde Q}$. Taking linear combinations of the lateral slices could, in principle, alter their rank; we will use the latter possibility to show that $\widetilde P$ must be of a constrained form.
Write $\widetilde P=\begin{bmatrix}P_{1,1}&P_{1,2}\\P_{2,1}&P_{2,2}\end{bmatrix}$ where $P_{1,1}$ is of size $n\times n$. We first claim that $P_{1,2}=0$. For if not, then in $(\mathtt A^{\mathrm{lat}})^{\widetilde P}$ (the side view), one of the last $n^2$ frontal slices receives a nonzero contribution from one of the first $n$ frontal slices of $\mathtt A^{\mathrm{lat}}$. Looking at the form of these slices from Equation3, we see that any such nonzero combination will have rank $\ge n$, but this is a contradiction since the corresponding slice in $\mathtt B^{\mathrm{lat}}$ has rank 1. Thus $P_{1,2}=0$, and therefore $P_{1,1}$ must be invertible, since $\widetilde P$ is.
Finally, we claim that $P_{1,1}$ has to be a monomial matrix. If not, then some frontal slice of $(\mathtt A^{\mathrm{lat}})^{\widetilde P}$ among the first $n$ would have a contribution from more than one of these $n$ slices. Considering the lower-right $n^2\times n^2$ sub-matrix of such a slice, we see that it would have rank exactly $kn$ for some $k\ge2$, which is again a contradiction since the first $n$ slices of $\mathtt B^{\mathrm{lat}}$ all have rank $<2n$. It follows that $P_{1,1}^tA_iP_{1,1}$, $i\in[m]$, are in $\mathcal B$, and thus $\mathcal A$ and $\mathcal B$ are monomially isometric via $P_{1,1}$.
\end{proof}
